% Part: computability
% Chapter: recursive-functions
% Section: primes

\documentclass[../../../include/open-logic-section]{subfiles}

\begin{document}

\olfileid{cmp}{rec}{pri}
\olsection{اولی اعداد}

محدود مسوّریت تے محدود کمترین قدر دی تلاش سانوں ایہہ وکھاون
لئی کافی اوزار دیندیاں نیں کہ قدرتی طور سامنے آون والے
فنکشن تے تعلق بنیادی بازگشتی نیں۔ مثال لئی ``$x$
، $y$ نوں پورا ونڈدا اے'' والا تعلق ویکھو، جیہنوں
$x \mid y$ لکھدے آں۔ $x \mid y$ اوہدوں درست اے
جدوں $y$،~$x$ دا صحیح عددی مضرب ہووے؛ غیر صفر
تقسیم کرن والے دی صورت وچ، ایہہ $y$ نوں~$x$ نال
بغیر باقی تقسیم کرن دے برابر اے۔ (جے ایہہ تعلق درست
نہ ہووے تے تقسیم کرن والا غیر صفر ہووے، یعنی $x$
نال $y$ دی تقسیم وچ باقی $> 0$ ہووے، تاں اسی
$x \nmid y$ لکھدے آں۔) دوجے لفظاں وچ، $x \mid y$
اوہدوں تے صرف اوہدوں اے جدوں کسے~$z$ لئی
$x \cdot z = y$ ہووے۔ جے ایہو جیہا شاہد موجود ہووے،
تاں اک شاہد~$z$ ایسا چنیا جا سکدا اے جیہڑا
$\leq y$ ہووے؛ صفر نال صفر دی صورت وچ صفر ای
چن لو۔ ایس لئی $x \mid y$ اوہدوں تے صرف اوہدوں
اے جدوں کسے $z \le y$ لئی $x \cdot z = y$
ہووے۔ اسی تعلق $x \mid y$ نوں $=$ تے ضرب توں
محدود وجودی مسوّریت راہیں ایویں متعین کر سکدے آں:
\[
x \mid y \defiff \bexists{z \leq y}{(x \cdot z) = y}.
\]
ایس طرح اسی وکھا دتا اے کہ $x \mid y$ بنیادی
بازگشتی اے۔

قدرتی عدد~$x$ نوں \emph{اولی} کہندے آں جے اوہ
نہ $0$ ہووے نہ $1$، تے اوہنوں صرف $1$ تے اوہ
آپ پورا ونڈن۔ دوجے لفظاں وچ، اولی عدد ایہو جیہے
نیں کہ جدوں وی $y \mid x$ ہووے، تاں یا $y = 1$
ہووے یا~$y=x$۔ ایہہ جانچن لئی کہ $x$~اولی اے
یا نہیں، سانوں صرف ایہہ جانچنا پیندا اے کہ
$y \mid x$ اے یا نہیں، اوہناں سارے $y \le x$
لئی؛ کیونکہ ایتھے حد
دے تحت جے $y > x$ ہووے، تاں آپے~$y \nmid x$
ہووے گا۔ ایس لئی تعلق $\fn{Prime}(x)$، جیہڑا
اوہدوں تے صرف اوہدوں درست اے جدوں $x$ اولی ہووے،
ایویں متعین ہو سکدا اے:
\[
\fn{Prime}(x) \defiff x \geq 2 \land \bforall{y \leq x}{(y \mid x \lif y
  = 1 \lor y = x)}
\]
تے ایس لئی ایہہ بنیادی بازگشتی اے۔

اولی اعداد $2$، $3$، $5$، $7$، $11$ وغیرہ نیں۔
فنکشن $p(x)$ ویکھو، جیہڑا ایس سلسلے دا $x$واں
اولی عدد واپس دیندا اے؛ یعنی $p(0) =
2$، $p(1) = 3$، $p(2) = 5$ وغیرہ۔ (سہولت
لئی اسی اکثر $p(x)$ نوں $p_x$ لکھاں گے
($p_0=2$، $p_1=3$ وغیرہ)۔)

جے ساڈے کول فنکشن
$\fn{nextPrime(x)}$ ہووے، جیہڑا~$x$ توں وڈا
پہلا اولی عدد واپس دیندا اے، تاں $p$ نوں بنیادی
بازگشت راہیں سوکھے نال متعین کیتا جا سکدا اے:
\begin{align*}
  p(0) & = 2\\
  p(x+1) & = \fn{nextPrime}(p(x))
\end{align*}
چونکہ $\fn{nextPrime}(x)$ اوہ سب توں چھوٹا
$y$ اے جیہڑا $y > x$ ہووے تے $y$~اولی ہووے،
ایہہ بغیر حد تلاش نال سوکھے حساب کیتا جا سکدا اے۔
پر اقلیدس دے اک نتیجے کر کے ایہہ محدود کمترین
قدر دی تلاش نال وی متعین ہو سکدا اے: $x$ توں
وڈا تے $\fact{x}+1$ توں ودھ نہ ہوون والا کوئی
اولی عدد ضرور ہوندا اے۔
\[
  \fn{nextPrime(x)} =
  \bmin{y \leq \fact{x}+1}{(y > x \land \fn{Prime}(y))}.
\]
ایہہ وکھاؤندا اے کہ $\fn{nextPrime}(x)$ تے ایس
لئی $p(x)$ صرف قابلِ حساب ای نہیں، بلکہ بنیادی
بازگشتی نیں۔

(جے دلچسپی ہووے، تاں اقلیدس دے نتیجے دا چھوٹا
ثبوت ایہہ اے۔ دو توں گھٹ حد دے دو کیس سیدھے
جانچے جا سکدے نیں؛ ہُن حد گھٹ توں گھٹ دو ہووے۔
فرض کرو $p_n$ سب توں وڈا اولی عدد اے جیہڑا
$\le x$ ہووے، تے ضرب حاصل $p =
p_0\cdot p_1 \cdot \dots \cdot p_n$ ویکھو؛ ایہہ
اوہناں سارے اولی اعداد~$\le x$ دا ضرب حاصل اے۔ یا تاں
$p+1$ خود اولی اے یا $x$ تے~$p+1$ دے وچ
کوئی اولی عدد اے۔ کیوں؟ فرض کرو $p+1$ اولی
نہیں۔ تاں کوئی اولی عدد $q \mid p+1$ اے جیہڑا
$q < p+1$ ہووے۔ $\le x$ والے اولی اعداد وچوں
کوئی وی $p+1$ نوں پورا نہیں ونڈدا۔ ($p$ دی تعریف
مطابق، ہر اولی عدد $p_i \le x$،~$p$ نوں پورا
ونڈدا اے، یعنی باقی~$0$ نال۔ ایس لئی ہر
$p_i \le x$، $p+1$ نوں تقسیم کرن اُتے باقی~$1$
دیندا اے، تے ایس کرکے $p_i \nmid p+1$۔) لہٰذا
$q$ اولی اے جیہڑا $> x$ تے $< p+1$ اے۔
نالے $p \le \fact{x}$، ایس لئی کوئی اولی عدد
$> x$ تے $\le
\fact{x}+1$ ضرور اے۔)

\begin{prob}
محدود کمترین قدر دی تلاش ورت کے صحیح عددی تقسیم
$d(x, y)$ متعین کرو۔
\end{prob}

\end{document}
